\section{\texorpdfstring{Tsallis 熵与 Rényi 熵}{Tsallis 熵与 Rényi 熵}}\label{app3-entropy-12-01532}
Rényi 的 \(\alpha\) 熵与 Tsallis 的 \(\alpha\) 熵通过以下单调关系相联系：

\setcounter{equation}{115}
\begin{equation}
\label{FD113-entropy-12-01532}
\fitmath{H_{R}^{(\alpha)}{(\overline{\mathbf{p}})} = \frac{1}{1 - \alpha}\log\left( {1 + {(1 - \alpha)}H_{T}^{(\alpha)}{(\overline{\mathbf{p}})}} \right)}
\end{equation}

或者使用 \(\alpha\) 对数写为

\setcounter{equation}{116}
\begin{equation}
\label{FD114-entropy-12-01532}
\fitmath{H_{T}^{(\alpha)}{(\overline{\mathbf{p}})} = \log_{\,\alpha}e^{\; H_{R}^{(\alpha)}{(\overline{\mathbf{p}})}}}
\end{equation}

\(\Psi_{R}^{(\alpha)} = - H_{R}^{(\alpha)}\) 是 \(\overline{\mathbf{p}}\) 的凸函数，因此可以构造相应的散度。


Bregman 散度可以定义为以下更一般的形式。


若 \(\left. \Phi{(\overline{\mathbf{p}})}:\;\mathbb{R}^{n}\rightarrow\mathbb{R} \right.\) 是严格凸的 \(C^{1}\) 函数，即一阶可微函数，则离散正测度上相应的 Bregman 距离或散度定义为

\setcounter{equation}{117}
\begin{equation}
\label{FD115-entropy-12-01532}
\fitmath{D_{\Phi}(\overline{\mathbf{p}}{\,||\,}\overline{\mathbf{q}}{) = \Phi}{(\overline{\mathbf{p}})} - \Phi{(\overline{\mathbf{q}})} - {(\overline{\mathbf{p}} - \overline{\mathbf{q}})}^{T}\nabla\Phi{(\overline{\mathbf{q}})}}
\end{equation}

其中 \(\nabla\Phi(\overline{\mathbf{q}})\) 表示 \(\Phi\) 在 \(\overline{\mathbf{q}}\) 处的梯度。


由 Rényi 熵导出的离散概率密度上的 Bregman 散度为

\setcounter{equation}{118}
\begin{equation}
\label{FD116-entropy-12-01532}
\fitmath{D_{R}^{(\alpha)}(\overline{\mathbf{p}}{\,||\,}\overline{\mathbf{q}}{) =}\frac{1}{1 - \alpha}\log\frac{h_{\alpha}{(\overline{\mathbf{q}})}}{h_{\alpha}{(\mathbf{p})}} - \frac{\alpha}{1 - \alpha}\frac{1}{h_{\alpha}{(\overline{\mathbf{q}})}}{(1 - \sum\limits_{i}{\overline{p}}_{i}\,{\overline{q}}_{i}^{\alpha - 1})}}
\end{equation}

Tsallis 熵与概率分布的 \(\alpha\) 指数族密切相关。由向量参数 \(\mathbf{\theta}\) 参数化的一族概率分布 \(S = \left\{ {\overline{p}{(\mathbf{x},\mathbf{\theta})}} \right\}\)，若能写为

\setcounter{equation}{119}
\begin{equation}
\label{FD117-entropy-12-01532}
\fitmath{\log_{\,\alpha}\overline{p}{(\mathbf{x},\mathbf{\theta})} = \sum\limits_{i}\theta_{i}x_{i} - \Psi{(\mathbf{\theta})}}
\end{equation}

其中 \(\Psi(\mathbf{\theta})\) 为凸函数，则称为 \(\alpha\) 指数族。这与文献[\hyperlink{B6-entropy-12-01532}{6}]定义的 \(\alpha\) 族相同。对于任意 \(\alpha\)，离散概率分布集合 \(S = \left\{ \mathbf{p} \right\}\) 都构成 \(\alpha\) 指数族。为说明这一点，引入随机变量 \(\mathbf{x} = (x)\)：

\setcounter{equation}{120}
\begin{equation}
\label{FD118-entropy-12-01532}
\fitmath{x_{i} = \delta_{i}}
\end{equation}

当结果为 i 时，它取值为 1，该结果的概率为 \(\overline{p}\)。


于是，概率可以写为

\setcounter{equation}{121}
\begin{equation}
\label{FD119-entropy-12-01532}
\fitmath{\log_{\,\alpha}\overline{p}{(\mathbf{x},\mathbf{\theta})} = \sum\limits_{i}\theta_{i}\delta_{i} - \Psi{(\mathbf{\theta})}}
\end{equation}

其中

\setcounter{equation}{122}
\begin{equation}
\label{FD120-entropy-12-01532}
\fitmath{\theta_{i} = \log_{\,\alpha}{\overline{p}}_{i}}
\end{equation}

由于 \(\sum{\overline{p}}_{i} = 1\) 成立，或者写成

\setcounter{equation}{123}
\begin{equation}
\label{FD121-entropy-12-01532}
\fitmath{\sum\delta_{i} = 1}
\end{equation}

因此，\(\mathbf{\theta}\) 的自由变量为 \(n - 1\)，记为 \(\left( {\theta_{1},\cdots,\theta_{n - 1}} \right)\)。此时，\(\theta_{n}\) 是其余变量的函数，得到

\setcounter{equation}{124}
\begin{equation}
\label{FD122-entropy-12-01532}
\fitmath{\Psi{(\mathbf{\theta})} = - \log_{\,\alpha}\theta_{n}}
\end{equation}




可以由这个凸函数 \(\Psi(\mathbf{\theta})\) 导出散度。不过，使用 \(\Psi(\mathbf{\theta})\) 的对偶函数更直观，该对偶函数为

\setcounter{equation}{125}
\begin{equation}
\label{FD123-entropy-12-01532}
\fitmath{\varphi{(\mathbf{\eta})} = \frac{1}{1 - \alpha}\left( {\frac{1}{h_{\alpha}{(\overline{\mathbf{p}})}} - 1} \right)}
\end{equation}

其中 \(\mathbf{\eta}\) 是 \(\mathbf{\theta}\) 通过 Legendre 变换得到的对偶变量：

\setcounter{equation}{126}
\begin{equation}
\label{FD124-entropy-12-01532}
\fitmath{\mathbf{\eta} = \nabla_{\mathbf{\theta}}\Psi{(\mathbf{\theta})}}
\end{equation}




得到的散度为

\setcounter{equation}{127}
\begin{equation}
\label{FD125-entropy-12-01532}
\fitmath{D_{\exp}^{(\alpha)}{(\overline{\mathbf{p}}\, \parallel \,\overline{\mathbf{q}})} = \frac{1}{h_{\alpha}{(\overline{\mathbf{p}})}}D_{T}^{(\alpha)}{(\overline{\mathbf{p}}\, \parallel \,\overline{\mathbf{q}})}}
\end{equation}

这是 \(\alpha\) 散度的共形变换[\hyperlink{B7-entropy-12-01532}{7}]。共形变换也称共形映射或保角变换，它保持任意两条直线或曲线在交点处的局部夹角。
